\documentclass[a4paper,14pt,oneside]{extarticle}

\usepackage[utf8]{inputenc}
\usepackage[T2A]{fontenc}
\usepackage[english,russian]{babel}
\usepackage[left=3cm,right=1.5cm,top=2cm,bottom=2cm]{geometry}
\usepackage{setspace}
\usepackage{amsmath}
\usepackage{graphicx}
\usepackage{float}
\usepackage{array}
\usepackage{longtable}
\usepackage{verbatim}
\usepackage{fvextra}
\usepackage[hidelinks]{hyperref}

\setstretch{1.35}
\setlength{\parindent}{1.25cm}
\setlength{\parskip}{0pt}
\addto\captionsrussian{
  \renewcommand{\contentsname}{Содержание}
  \renewcommand{\figurename}{Рисунок}
}

\begin{document}
\selectlanguage{russian}

\begin{titlepage}
\begin{center}
\textbf{МИНИСТЕРСТВО НАУКИ И ВЫСШЕГО ОБРАЗОВАНИЯ РФ}\\
\vspace{0.4cm}
\textbf{[Наименование образовательной организации]}\\
\textbf{[Наименование кафедры]}

\vfill

\textbf{\Large ОТЧЁТ}\\
\vspace{0.4cm}
\textbf{\large по лабораторной работе № 2}\\
\vspace{0.4cm}
\textbf{\large «Простейшие криптографические преобразования»}

\vfill

\begin{flushright}
\begin{minipage}{0.48\textwidth}
Выполнил: Сечейко Н.В.\\
\vspace{0.3cm}
Проверил(а): \underline{\hspace{4cm}}
\end{minipage}
\end{flushright}

\vfill

\textbf{Минск}\\
\textbf{2026}
\end{center}
\end{titlepage}

\tableofcontents
\newpage

\section{Цель и содержание работы}

Цель работы --- изучить простейшие криптографические преобразования и
разработать программу для шифрования, расшифрования и атаки шифра Виженера.

В соответствии с методическими указаниями программа должна:
\begin{enumerate}
\item реализовывать шифрование и расшифрование текста;
\item реализовывать атаку полным перебором ключа;
\item позволять оценить криптографическую стойкость шифра;
\item содержать алгоритм возможного усложнения шифра.
\end{enumerate}

\section{Теоретические сведения}

Шифр Виженера является полиалфавитным шифром подстановки. В отличие от
шифра Цезаря, в нём используется несколько алфавитов подстановки, которые
циклически выбираются в соответствии с символами ключевого слова.

Каждому символу алфавита ставится в соответствие число от 0 до $n-1$, где
$n$ --- мощность алфавита. В работе используется английский алфавит в нижнем
регистре:

\[
\texttt{abcdefghijklmnopqrstuvwxyz}.
\]

Мощность алфавита равна:

\[
n=26.
\]

Пусть $x_i$ --- числовое значение символа открытого текста, $k_i$ ---
числовое значение соответствующего символа циклически повторяемого ключа, а
$y_i$ --- числовое значение символа шифртекста. Тогда шифрование выполняется
по формуле:

\[
y_i=(x_i+k_i)\bmod n.
\]

Расшифрование выполняется по формуле:

\[
x_i=(y_i-k_i)\bmod n.
\]

Если длина ключа меньше длины сообщения, ключ повторяется. Например, для
текста \texttt{attackatdawn} и ключа \texttt{dog} используется
последовательность:

\[
\texttt{dogdogdogdog}.
\]

Основная уязвимость шифра Виженера заключается в циклическом повторении
ключа. При известной длине ключа возможен полный перебор всех вариантов
ключа. Если длина ключа равна $m$, количество ключей составляет:

\[
N_{\text{ключей}}=n^m.
\]

Среднее число попыток при полном переборе:

\[
N_{\text{ср}}=\frac{n^m}{2}.
\]

\section{Исходные данные и вариант}

Вариант 26 --- шифр Виженера.

В качестве языка исходного текста выбран английский язык. Программа работает
со строчными буквами английского алфавита без пробелов и знаков препинания.

\section{Описание разработанной программы}

Программа написана на языке C++. После запуска пользователю доступно меню:

\begin{enumerate}
\item шифрование открытого текста;
\item расшифрование шифртекста;
\item атака полным перебором ключа;
\item завершение работы.
\end{enumerate}

При шифровании программа получает открытый текст и ключ. Для каждого символа
вычисляется сумма его номера и номера соответствующего символа ключа по
модулю 26.

При расшифровании программа получает шифртекст и ключ. Из номера символа
шифртекста вычитается номер символа ключа по модулю 26.

Для атаки полным перебором программа получает шифртекст, известный открытый
текст и длину ключа. Затем она последовательно формирует все ключи заданной
длины от символов \texttt{a} до \texttt{z}. Для каждого ключа выполняется
расшифрование, после чего результат сравнивается с известным открытым
текстом.

\section{Ход выполнения и результаты}

Для проверки программы использованы следующие исходные данные:

\begin{itemize}
\item открытый текст: \texttt{attackatdawn};
\item ключ: \texttt{dog};
\item длина ключа: $m=3$.
\end{itemize}

Результат шифрования:

\[
\texttt{attackatdawn}\longrightarrow\texttt{dhzdqqdhjdkt}.
\]

Проверка расшифрования:

\[
\texttt{dhzdqqdhjdkt}\longrightarrow\texttt{attackatdawn}.
\]

При атаке полным перебором в качестве исходного открытого текста указано
\texttt{attackatdawn}. Программа находит ключ \texttt{dog}. Полное количество
проверяемых ключей равно:

\[
26^3=17576.
\]

Среднее количество попыток:

\[
\frac{26^3}{2}=8788.
\]

Ключ \texttt{dog} находится в переборе после 2399 проверок, если ключи
перебирать в лексикографическом порядке. В программе после обнаружения
совпадения перебор продолжается до проверки всех ключей заданной длины.

\section{Оценка криптографической стойкости}

Стойкость реализованного шифра зависит от длины ключа. Для английского
алфавита количество возможных ключей определяется выражением $26^m$.

Для ключа длиной 1:

\[
26^1=26.
\]

Для ключа длиной 3:

\[
26^3=17576.
\]

Для ключа длиной 5:

\[
26^5=11881376.
\]

Таким образом, короткий ключ шифра Виженера может быть найден полным
перебором. Дополнительную информацию атакующему даёт повторение ключа в
шифртексте. Для более длинных сообщений возможно определить длину ключа
методом прореживания и затем провести криптоанализ совокупности шифров
Цезаря.

Реализованный шифр не предназначен для защиты информации в современных
информационных системах. Его криптографическая стойкость существенно ниже
стойкости современных симметричных алгоритмов.

\section{Алгоритм усложнения шифра}

Для усложнения базового шифра Виженера можно использовать ключ, длина
которого равна длине сообщения. Алгоритм имеет следующий вид:

\begin{enumerate}
\item сгенерировать случайную последовательность символов длиной не меньше
длины открытого текста;
\item использовать эту последовательность как ключ;
\item для каждого символа текста выполнить сложение по модулю 26;
\item передать получателю шифртекст и ключ по разным защищённым каналам;
\item после однократного использования уничтожить ключ.
\end{enumerate}

При условии случайности ключа, равенства его длины длине сообщения и
однократного использования такой подход соответствует идее одноразового
шифроблокнота и существенно усложняет атаку по повторяющемуся ключу.

\section{Заключение}

В ходе лабораторной работы для варианта 26 реализован шифр Виженера на
английском алфавите. Программа выполняет шифрование и расшифрование текста,
а также атаку полным перебором ключа по известному открытому тексту.

На примере текста \texttt{attackatdawn} и ключа \texttt{dog} получен
шифртекст:

\[
\texttt{dhzdqqdhjdkt}.
\]

Полученный текст успешно расшифровывается исходным ключом. Для ключа длиной 3
существует 17576 вариантов, поэтому такой ключ может быть найден полным
перебором.

\appendix
\section{Исходный код программы}

\VerbatimInput[breaklines=true,breakanywhere=true]{lab2.cpp}

\end{document}
